% 第3章：本章小结。
\FloatBarrier
\section{本章小结}
\label{sec:graph-opt-sp-summary}
\label{sec:shortest-paths-summary}
\label{sec:apsp-summary}

本章围绕有向加权图上的最短路问题，先后解决了从单一源点出发的单源最短路问题，以及同时求所有顶点对距离的全源最短路问题。

在共同基础（第~\ref{sec:shortest-path-foundations}节）中，按有限行走成本的下确界定义最短路距离后，每个有序顶点对$(s,v)$的距离属于三类之一：不可达时为$+\infty$；可达且不受负环影响时为有限值，并由一条简单路径达到；能从某个$s$可达的负环到达时为$-\infty$。负环只影响能从源点经该环到达的终点。各类算法共用的核心机制是距离标签与逐弧松弛：标签从“源为零、其余为无穷”出发，只通过松弛下降，因而始终是有实际路线见证的距离上界；而当所有弧都满足三角不等式$d(v)\leq d(u)+\ell_{u,v}$时，沿任意行走裂项相加又给出距离下界，两者合一即构成不依赖算法过程的最优性证书。

在单源最短路（第~\ref{sec:single-source-shortest-paths}节）中，三种算法的差别在于何时能认定标签不再改变，而这取决于输入图的环结构与边长符号。有向无环图按拓扑序处理，前驱的距离在使用前已经确定，弧长可取任意实数，总时间$O(n+a)$，编辑距离一类的状态递推正是这种情形。

Bellman--Ford算法不要求无环，同步版本第$k$轮的标签恰是至多$k$条弧的最优值，$n-1$轮后仍能改善的弧暴露负环的影响，就地更新版本总时间$O(n+na)$，直接复制数组的同步版本为$O(n(n+a))$。

Dijkstra算法要求可达部分弧长非负，每次出队的最小标签顶点立即永久确定，用二叉堆时为$O((n+a)\log n)$；哪怕没有负环，一条负边就可能使它过早锁定标签而出错。

在全源最短路（第~\ref{sec:all-pairs-shortest-paths}节）中，为了突破含负边图上重复运行$n$次Bellman--Ford算法的$O(n^2+n^2m)$时间上界，本章建立了两条互补路线。

其一是面向稠密图与矩阵计算的Floyd--Warshall算法：它按允许经过的中间顶点编号集合$\{1,\ldots,k\}$逐层递推，利用简单路径至多经过枢轴顶点$k$一次的性质做二选一更新，在代数上对应$(\min,+)$热带半环上的矩阵消元；就地更新须把$k$放在最外层，$i,j$两层可以互换，由对角线$d(i,i)<0$检测负环，最坏时间为$\Theta(n^3)$、空间为$\Theta(n^2)$。

其二是面向含负边稀疏图的Johnson算法：它引入顶点势$h(v)$定义约化权$\tilde w(u,v)=w(u,v)+h(u)-h(v)$，依靠路径裂项相消保持最短路与环权不变；通过添加超级源点$s_0$运行一次Bellman--Ford取$h(v)=d^*_{\mathcal G'}(s_0,v)$，由三角不等式保证$\tilde w(u,v)\geq0$，随后调用$n$次Dijkstra算法将总时间降至$O(nm+n^2\log n)$。

下一章（第~\ref{chap:minimum-spanning-tree}章）转向另一种选择对象：最小生成树比较的是一整组连通边的总造价，而不是两点间路线的长度。本章建立的逐弧三角不等式与Johnson势函数重赋权，则会在第~\ref{chap:bipartite-matching-assignment}章的可行标号和第~\ref{chap:min-cost-flow}章的节点势与约化费用中再次发挥核心作用。
